\mxfront{stmg-bank-c}

% -----------------------------------------------------------------------------
% Chapitre 5 -- Statistiques à deux variables quantitatives
% -----------------------------------------------------------------------------
\mxexdef{31}{Construire un nuage de points}{%
Une supérette compare, sur six mois, l'écart $x_i$ entre la température moyenne du mois et la normale saisonnière (en °C) et la variation $y_i$ de ses ventes de boissons fraîches par rapport au même mois de l'année précédente (en centaines de bouteilles).
\[
\begin{array}{c|rrrrrr}
\text{Écart de température }x_i\text{ (°C)}&-4&-1&0&2&5&7\\\hline
\text{Variation des ventes }y_i\text{ (centaines)}&-3&-2&0&2&4&5
\end{array}
\]
\begin{enumerate}
\item Indiquer la variable placée en abscisse, celle placée en ordonnée, ainsi que leurs unités.
\item Choisir des échelles adaptées puis représenter le nuage de points $M_i\pt{x_i;y_i}$ dans un repère.
\item Calculer les coordonnées du point moyen $G$ du nuage et le placer.
\item Décrire en une phrase l'allure du nuage.
\end{enumerate}
}{%
\begin{enumerate}
\item En abscisse : l'écart de température $x$, en °C ; en ordonnée : la variation des ventes $y$, en centaines de bouteilles. Les deux variables prennent des valeurs négatives : les axes doivent se couper à l'intérieur du graphique.
\item On peut prendre $1$~cm pour $1$~°C (de $-5$ à $8$) et $1$~cm pour une centaine de bouteilles (de $-4$ à $6$). On place les six points $\pt{-4;-3}$, $\pt{-1;-2}$, $\pt{0;0}$, $\pt{2;2}$, $\pt{5;4}$ et $\pt{7;5}$, sans les relier.
\begin{center}
\begin{tikzpicture}
\begin{axis}[width=.62\textwidth,height=5.2cm,axis lines=middle,xmin=-5,xmax=8,ymin=-4,ymax=6,grid=major,xtick={-4,...,7},ytick={-3,...,5},xlabel={$x$},ylabel={$y$}]
\addplot[only marks,mark=*,mxteacher,mark size=2.2pt] coordinates {(-4,-3)(-1,-2)(0,0)(2,2)(5,4)(7,5)};
\addplot[only marks,mark=square*,mxwarn,mark size=2.6pt] coordinates {(1.5,1)};
\node[mxwarn,anchor=south east] at (axis cs:1.5,1) {$G$};
\end{axis}
\end{tikzpicture}
\end{center}
\item $\bar x=\dfrac{-4-1+0+2+5+7}{6}=\dfrac{9}{6}=\num{1,5}$ et $\bar y=\dfrac{-3-2+0+2+4+5}{6}=\dfrac{6}{6}=1$. Donc $G\pt{\num{1,5};1}$.
\item Le nuage est orienté vers le haut : plus le mois est chaud par rapport à la normale, plus les ventes de boissons fraîches augmentent ; les points sont proches d'une droite.
\end{enumerate}
}{5}

\mxexdef{32}{Interpréter un nuage sans conclure à une causalité}{%
Pendant huit semaines d'été, un glacier de bord de mer relève ses ventes hebdomadaires $x$ (en centaines de cornets) ; la pharmacie voisine relève le nombre $y$ de crèmes apaisantes pour coups de soleil vendues la même semaine. Le nuage de points obtenu est représenté ci-dessous.
\begin{center}
\begin{tikzpicture}
\begin{axis}[width=.7\textwidth,height=5.4cm,axis lines=left,xmin=10,xmax=31,ymin=0,ymax=24,grid=major,xtick={10,12,...,30},xlabel={ventes de glaces (centaines)},ylabel={crèmes vendues}]
\addplot[only marks,mark=*,mxteacher,mark size=2.4pt] coordinates {(12,8)(15,10)(18,13)(20,13)(23,16)(26,18)(28,21)(29,9)};
\end{axis}
\end{tikzpicture}
\end{center}
\begin{enumerate}
\item Décrire le nuage : sens de variation, forme et dispersion.
\item Un point s'écarte nettement des autres. Donner ses coordonnées et proposer une explication plausible.
\item Le pharmacien affirme : « Ce sont les glaces qui provoquent les coups de soleil. » Que penser de cette affirmation ?
\item Proposer un facteur qui pourrait expliquer l'évolution simultanée des deux séries.
\end{enumerate}
}{%
\begin{enumerate}
\item Hormis un point, le nuage est croissant : les semaines où le glacier vend davantage sont aussi celles où la pharmacie vend plus de crèmes. Les points sont assez proches d'une droite, avec une dispersion faible.
\item Le point $\pt{29;9}$ est atypique : très fortes ventes de glaces mais peu de crèmes vendues. Explication possible : une semaine nuageuse marquée par un festival qui a attiré beaucoup de promeneurs sans exposition forte au soleil.
\item Le nuage montre seulement que les deux variables varient dans le même sens (corrélation). Il ne prouve pas qu'une variable soit la cause de l'autre : l'affirmation n'est pas justifiée par ces données.
\item L'ensoleillement (ou la chaleur, la fréquentation touristique) fait augmenter à la fois les ventes de glaces et les coups de soleil : c'est un facteur commun aux deux séries.
\end{enumerate}
}{4}

\mxexdef{33}{Juger une droite d'ajustement et l'exploiter}{%
Le gérant d'une boutique chauffée au gaz relève, certains jours d'hiver, la température extérieure moyenne $x_i$ (en °C) et la consommation de gaz $y_i$ de la boutique (en kWh).
\[
\begin{array}{c|rrrrrr}
x_i\text{ (°C)}&-4&-2&0&3&6&9\\\hline
y_i\text{ (kWh)}&62&55&50&41&32&24
\end{array}
\]
\begin{enumerate}
\item Représenter le nuage de points dans un repère, puis calculer les coordonnées du point moyen $G$.
\item Un stagiaire propose la droite $d_1$ d'équation $y=-2x+50$. La tracer. Cet ajustement paraît-il satisfaisant ? Justifier en utilisant le point $G$.
\item On retient désormais la droite $d_2$ d'équation $y=-3x+50$. Vérifier que $G$ appartient à $d_2$ et interpréter les nombres $-3$ et $50$ dans le contexte.
\item À l'aide de $d_2$, estimer :
  \begin{itemize}
  \item la consommation un jour où il fait $1$~°C ;
  \item la consommation un jour où il fait $12$~°C ;
  \item la température d'un jour où la consommation est de $35$~kWh.
  \end{itemize}
  Préciser, pour les deux premières estimations, s'il s'agit d'une interpolation ou d'une extrapolation.
\end{enumerate}
}{%
\begin{enumerate}
\item On place les six points. $\bar x=\dfrac{-4-2+0+3+6+9}{6}=\dfrac{12}{6}=2$ et $\bar y=\dfrac{62+55+50+41+32+24}{6}=\dfrac{264}{6}=44$, donc $G\pt{2;44}$. Le nuage est décroissant et presque aligné.
\item Pour $x=2$, $d_1$ donne $y=-2\times2+50=46\neq44$ : $d_1$ ne passe pas par $G$. De plus, elle est trop peu inclinée : elle passe au-dessous des premiers points ($58$ pour $62$ à $-4$~°C) et nettement au-dessus des derniers ($32$ pour $24$ à $9$~°C). L'ajustement n'est pas satisfaisant.
\item $-3\times2+50=44$ : $G\in d_2$. Le coefficient $-3$ signifie qu'un degré de plus correspond, en moyenne, à environ $3$~kWh consommés en moins ; $50$ est la consommation estimée un jour à $0$~°C.
\item Pour $x=1$ : $y=-3+50=47$, soit environ $47$~kWh (interpolation, car $1\in\intff{-4}{9}$).\\
Pour $x=12$ : $y=-36+50=14$, soit environ $14$~kWh (extrapolation, car $12>9$ : à interpréter avec prudence).\\
$-3x+50=35$ équivaut à $x=5$ : une consommation de $35$~kWh correspond à environ $5$~°C.
\end{enumerate}
}{6}

\mxexdef{34}{Régression affine et prévision chronologique}{%
Un opérateur régional relève le nombre d'abonnés à son offre fibre optique au 31 décembre de chaque année.
\[
\begin{array}{c|cccccc}
\text{Année}&2018&2019&2020&2021&2022&2023\\\hline
\text{Rang de l'année }x_i&0&1&2&3&4&5\\
\text{Abonnés }y_i\text{ (en milliers)}&\num{8,2}&\num{11,5}&\num{13,9}&\num{17,4}&\num{20,1}&\num{23,8}
\end{array}
\]
\begin{enumerate}
\item Représenter le nuage de points $M_i\pt{x_i;y_i}$. Un ajustement affine est-il envisageable ?
\item À l'aide de la calculatrice, déterminer une équation de la droite d'ajustement de $y$ en $x$ obtenue par la méthode des moindres carrés. Arrondir les coefficients à $10^{-2}$.
\item En supposant que la tendance se maintient, estimer le nombre d'abonnés fin 2026.
\item Selon ce modèle, en quelle année le nombre d'abonnés dépasserait-il $\num{40000}$ ?
\end{enumerate}
}{%
\begin{enumerate}
\item Les six points sont presque alignés selon une droite croissante : un ajustement affine est pertinent.
\item La calculatrice donne $a\approx\num{3,0657}$ et $b\approx\num{8,1524}$, d'où la droite $y=\num{3,07}x+\num{8,15}$.
\item Fin 2026 correspond au rang $x=8$ : $y=\num{3,07}\times8+\num{8,15}=\num{32,71}$. On peut prévoir environ $\num{32710}$ abonnés. C'est une extrapolation : le résultat suppose que la tendance observée de 2018 à 2023 se poursuit.
\item On résout $\num{3,07}x+\num{8,15}>40$, soit $x>\dfrac{40-\num{8,15}}{\num{3,07}}\approx\num{10,37}$. Le premier rang entier convenable est $x=11$ (contrôle : $\num{3,07}\times10+\num{8,15}=\num{38,85}<40$ et $\num{3,07}\times11+\num{8,15}=\num{41,92}>40$). Le seuil serait dépassé en $2018+11=2029$.
\end{enumerate}
}{6}

\mxexdef{35}{Droite par deux points et extrapolation prudente}{%
Une plateforme logistique a relevé le nombre de colis expédiés chaque année (en milliers).
\[
\begin{array}{c|ccccc}
\text{Année}&2019&2020&2021&2022&2023\\\hline
\text{Rang de l'année }x_i&0&1&2&3&4\\
\text{Colis expédiés }y_i\text{ (milliers)}&\num{212,4}&\num{268,0}&\num{341,7}&\num{398,2}&\num{463,6}
\end{array}
\]
\begin{enumerate}
\item Représenter le nuage de points $M_1,\ldots,M_5$ dans un repère orthogonal.
\item On ajuste ce nuage par la droite $(M_1M_5)$ passant par le premier et le dernier point. Déterminer une équation de cette droite (coefficients arrondis au dixième si nécessaire).
\item Avec ce modèle, estimer le nombre de colis expédiés en 2027.
\item Selon ce modèle, au cours de quelle année le nombre de colis expédiés dépasserait-il un million ?
\item Donner deux raisons pour lesquelles ces prévisions peuvent se révéler peu fiables.
\end{enumerate}
}{%
\begin{enumerate}
\item Les cinq points sont presque alignés selon une droite croissante.
\item $M_1\pt{0;\num{212,4}}$ et $M_5\pt{4;\num{463,6}}$. Le coefficient directeur vaut $\dfrac{\num{463,6}-\num{212,4}}{4-0}=\dfrac{\num{251,2}}{4}=\num{62,8}$ et l'ordonnée à l'origine est $\num{212,4}$ (abscisse $0$). Équation : $y=\num{62,8}x+\num{212,4}$.
\item 2027 correspond au rang $x=8$ : $y=\num{62,8}\times8+\num{212,4}=\num{714,8}$, soit environ $\num{714800}$ colis.
\item Un million de colis correspond à $y=\num{1000}$. On résout $\num{62,8}x+\num{212,4}>\num{1000}$, soit $x>\dfrac{\num{787,6}}{\num{62,8}}\approx\num{12,54}$. Le premier rang entier est $x=13$ (contrôle : $\num{62,8}\times12+\num{212,4}=966<\num{1000}$ et $\num{62,8}\times13+\num{212,4}=\num{1028,8}$). Le seuil serait franchi en $2019+13=2032$.
\item Ce sont des extrapolations loin de la période observée (2019--2023) : la tendance peut changer (concurrence, crise économique, saturation de la capacité de la plateforme) ; de plus la droite est construite à partir de deux points seulement.
\end{enumerate}
}{6}

\mxexdef{36}{Linéariser par changement de variable logarithmique}{%
Une agence de communication étudie, pour plusieurs clients, le budget mensuel de publicité en ligne $x_i$ (en centaines d'euros) et le nombre mensuel de visites $y_i$ du site du client (en milliers).
\[
\begin{array}{c|rrrrrrrrrr}
x_i&1&2&4&5&8&10&15&20&30&40\\\hline
y_i&13&19&28&29&35&36&42&44&49&52
\end{array}
\]
Les visites augmentent de moins en moins vite quand le budget augmente : un ajustement affine de $y$ en $x$ n'est pas pertinent. On pose $u_i=\log(x_i)$.
\begin{enumerate}
\item Calculer les valeurs $u_i$, arrondies au centième.
\item Représenter le nuage de points $\pt{u_i;y_i}$. À l'aide de la calculatrice, déterminer une équation de la droite d'ajustement de $y$ en $u$ par la méthode des moindres carrés (coefficients arrondis au dixième).
\item En déduire une expression de $y$ en fonction de $x$.
\item Estimer le nombre de visites pour un budget de $\num{1200}$~€, puis de $\num{6000}$~€. Commenter la fiabilité de chaque estimation.
\end{enumerate}
}{%
\begin{enumerate}
\item $u_i$ : $0$ ; $\num{0,30}$ ; $\num{0,60}$ ; $\num{0,70}$ ; $\num{0,90}$ ; $1$ ; $\num{1,18}$ ; $\num{1,30}$ ; $\num{1,48}$ ; $\num{1,60}$.
\item Le nuage $\pt{u_i;y_i}$ est presque aligné. La calculatrice donne (avec les $u_i$ arrondis) $a\approx\num{24,59}$ et $b\approx\num{12,42}$, d'où $y=\num{24,6}u+\num{12,4}$.
\item Comme $u=\log(x)$, on obtient $y=\num{24,6}\log(x)+\num{12,4}$.
\item $\num{1200}$~€ correspond à $x=12$ : $y=\num{24,6}\log(12)+\num{12,4}\approx\num{38,9}$, soit environ $\num{38900}$ visites (interpolation, $12\in\intff{1}{40}$).\\
$\num{6000}$~€ correspond à $x=60$ : $y=\num{24,6}\log(60)+\num{12,4}\approx\num{56,1}$, soit environ $\num{56100}$ visites. C'est une extrapolation ($60>40$) : estimation à interpréter avec prudence.
\end{enumerate}
}{6}

\mxexdef{37}{Étalonnage : ajustement, résidu et domaine de validité}{%
Une jeune entreprise étalonne un capteur de dioxyde de carbone (CO$_2$) à bas coût destiné aux bureaux. Placé dans des enceintes de concentration connue $C$ (en ppm, parties par million), le capteur délivre une tension $U$ (en volts).
\[
\begin{array}{c|cccccc}
\text{Tension }U\text{ (V)}&\num{0,42}&\num{0,81}&\num{1,23}&\num{1,60}&\num{2,05}&\num{2,41}\\\hline
\text{Concentration }C\text{ (ppm)}&400&800&\num{1200}&\num{1600}&\num{2000}&\num{2400}
\end{array}
\]
\begin{enumerate}
\item Représenter le nuage de points $\pt{U;C}$, avec $U$ en abscisse.
\item Un ajustement affine semble-t-il pertinent ? Déterminer à la calculatrice une équation de la droite d'ajustement de $C$ en $U$, sous la forme $C=aU+b$ (coefficients arrondis à l'unité).
\item Calculer le résidu pour $U=\num{1,23}$, c'est-à-dire l'écart entre la concentration réelle et celle donnée par le modèle. Commenter.
\item Dans un bureau, le capteur affiche $\num{1,4}$~V. Estimer la concentration en CO$_2$. Un seuil de $\num{1000}$~ppm est souvent retenu comme repère pour aérer : faut-il aérer ?
\item Peut-on utiliser ce modèle avec la même confiance pour une tension de $\num{3,5}$~V ? Justifier.
\end{enumerate}
}{%
\begin{enumerate}
\item Les six points sont presque alignés selon une droite croissante.
\item Oui. La calculatrice donne $a\approx\num{996,59}$ et $b\approx-\num{15,16}$, d'où $C=997U-15$.
\item Le modèle donne $997\times\num{1,23}-15=\num{1211,31}$. Le résidu vaut $\num{1200}-\num{1211,31}\approx-11$~ppm, soit moins de $\pc{1}$ de la valeur : l'ajustement est de bonne qualité.
\item $C=997\times\num{1,4}-15=\num{1380,8}$, soit environ $\num{1381}$~ppm (interpolation, car $\num{1,4}\in\intff{\num{0,42}}{\num{2,41}}$). Cette valeur dépasse $\num{1000}$~ppm : il faut aérer.
\item Le modèle donne $997\times\num{3,5}-15=\num{3474,5}$~ppm, mais $\num{3,5}$~V est hors de l'intervalle d'étalonnage $\intff{\num{0,42}}{\num{2,41}}$ : c'est une extrapolation. Rien ne garantit que le capteur reste linéaire à forte concentration ; il faudrait compléter l'étalonnage.
\end{enumerate}
}{6}

\mxexdef{38}{Point moyen et droite de régression : choix justifiés}{%
Pour chaque question, une seule proposition est exacte. La choisir et justifier.

Trois magasins d'une même enseigne ont une surface de vente $x$ (en centaines de m$^2$) et un chiffre d'affaires mensuel $y$ (en dizaines de milliers d'euros) représentés par les points $M_1\pt{2;5}$, $M_2\pt{4;11}$ et $M_3\pt{6;11}$.
\begin{enumerate}
\item Le point moyen $G$ de ce nuage est :
  \begin{enumerate}[label=\alph*.]
  \item le point $M_2$ ;
  \item le point de coordonnées $\pt{4;9}$ ;
  \item le point de coordonnées $\pt{12;27}$.
  \end{enumerate}
\item Une équation de la droite $(M_1M_3)$ est :
  \begin{enumerate}[label=\alph*.]
  \item $y=\num{1,5}x+2$ ;
  \item $y=2x+1$ ;
  \item $y=\num{1,5}x+3$.
  \end{enumerate}
\item La droite d'ajustement obtenue à la calculatrice par la méthode des moindres carrés :
  \begin{enumerate}[label=\alph*.]
  \item passe par le point $M_2$ ;
  \item passe par $G$ et a le même coefficient directeur que $(M_1M_3)$ ;
  \item passe par l'origine du repère.
  \end{enumerate}
\item Avec cette droite d'ajustement, le chiffre d'affaires estimé d'un magasin de $500$~m$^2$ est :
  \begin{enumerate}[label=\alph*.]
  \item $\num{95000}$~€ ;
  \item $\num{105000}$~€ ;
  \item $\num{110000}$~€.
  \end{enumerate}
\end{enumerate}
}{%
\begin{enumerate}
\item Réponse b : $\bar x=\dfrac{2+4+6}{3}=4$ et $\bar y=\dfrac{5+11+11}{3}=9$. La proposition c oublie de diviser les sommes par $3$.
\item Réponse a : le coefficient directeur vaut $\dfrac{11-5}{6-2}=\num{1,5}$ et $5=\num{1,5}\times2+b$ donne $b=2$. La droite $y=2x+1$ passe par $M_1$ mais pas par $M_3$ ($2\times6+1=13\neq11$).
\item Réponse b : la calculatrice donne $y=\num{1,5}x+3$. Elle passe par $G$ ($\num{1,5}\times4+3=9$), a pour coefficient directeur $\num{1,5}$ comme $(M_1M_3)$, et ne passe ni par $M_2$ ($\num{1,5}\times4+3=9\neq11$) ni par l'origine ($b=3\neq0$).
\item Réponse b : $500$~m$^2$ correspond à $x=5$, et $y=\num{1,5}\times5+3=\num{10,5}$ dizaines de milliers d'euros, soit $\num{105000}$~€ (interpolation).
\end{enumerate}
}{3}

% -----------------------------------------------------------------------------
% Chapitre 6 -- Probabilités conditionnelles
% -----------------------------------------------------------------------------
\mxexdef{39}{Arbre à trois branches, probabilité totale et inverse}{%
Un site de vente en ligne confie ses colis à trois transporteurs : $\pc{45}$ au transporteur $T_1$, $\pc{35}$ au transporteur $T_2$ et les autres au transporteur $T_3$. Le service client constate que $\pc{2}$ des colis acheminés par $T_1$ arrivent endommagés, contre $\pc{5}$ pour $T_2$ et $\pc{8}$ pour $T_3$.

On choisit un colis au hasard. On note $T_i$ l'événement « Le colis est acheminé par le transporteur $T_i$ » et $E$ l'événement « Le colis arrive endommagé ».
\begin{center}
\begin{tikzpicture}[grow'=right,level distance=30mm,level 1/.style={sibling distance=20mm},level 2/.style={sibling distance=9mm},every node/.style={font=\small},edge from parent/.style={draw,-{Stealth}}]
\node {}
 child {node {$T_1$}
   child {node {$E$} edge from parent node[above] {$\num{0,02}$}}
   child {node {$\overline E$} edge from parent node[below] {$\dots$}}
   edge from parent node[above] {$\num{0,45}$}}
 child {node {$T_2$}
   child {node {$E$} edge from parent node[above] {$\dots$}}
   child {node {$\overline E$} edge from parent node[below] {$\dots$}}
   edge from parent node[above] {$\dots$}}
 child {node {$T_3$}
   child {node {$E$} edge from parent node[above] {$\dots$}}
   child {node {$\overline E$} edge from parent node[below] {$\dots$}}
   edge from parent node[below] {$\dots$}};
\end{tikzpicture}
\end{center}
\begin{enumerate}
\item Recopier et compléter l'arbre pondéré ci-dessus.
\item Calculer la probabilité que le colis soit acheminé par $T_3$ et arrive endommagé.
\item Montrer que la probabilité qu'un colis arrive endommagé est égale à $\num{0,0425}$.
\item Un client signale un colis endommagé. Quelle est la probabilité qu'il ait été acheminé par $T_3$ ? Arrondir à $10^{-3}$.
\end{enumerate}
}{%
\begin{enumerate}
\item $P(T_3)=1-\num{0,45}-\num{0,35}=\num{0,20}$. Branches de second niveau : $P_{T_1}(\overline E)=\num{0,98}$ ; $P_{T_2}(E)=\num{0,05}$ et $P_{T_2}(\overline E)=\num{0,95}$ ; $P_{T_3}(E)=\num{0,08}$ et $P_{T_3}(\overline E)=\num{0,92}$.
\item $P(T_3\cap E)=P(T_3)\times P_{T_3}(E)=\num{0,20}\times\num{0,08}=\num{0,016}$.
\item $T_1$, $T_2$, $T_3$ forment une partition de l'univers. Par la formule des probabilités totales :
$P(E)=\num{0,45}\times\num{0,02}+\num{0,35}\times\num{0,05}+\num{0,20}\times\num{0,08}=\num{0,009}+\num{0,0175}+\num{0,016}=\num{0,0425}$.
\item $P_E(T_3)=\dfrac{P(T_3\cap E)}{P(E)}=\dfrac{\num{0,016}}{\num{0,0425}}\approx\num{0,376}$. Bien que $T_3$ ne transporte que $\pc{20}$ des colis, il est concerné par environ $\pc{37,6}$ des colis endommagés.
\end{enumerate}
}{5}

\mxexdef{40}{Compléter un arbre et calculer des intersections}{%
Dans cet exercice, les résultats sont donnés sous forme de fractions irréductibles. On considère l'arbre pondéré ci-dessous, où $E$ est un événement et où $R$, $S$ et $T$ forment une partition de l'univers.
\begin{center}
\begin{tikzpicture}[grow'=right,level distance=30mm,level 1/.style={sibling distance=30mm},level 2/.style={sibling distance=10mm},every node/.style={font=\small},edge from parent/.style={draw,-{Stealth}}]
\node {}
 child {node {$E$}
   child {node {$R$} edge from parent node[above] {$\tfrac14$}}
   child {node {$S$} edge from parent node[above] {$\dots$}}
   child {node {$T$} edge from parent node[below] {$\tfrac13$}}
   edge from parent node[above] {$\tfrac49$}}
 child {node {$\overline E$}
   child {node {$R$} edge from parent node[above] {$\tfrac3{10}$}}
   child {node {$S$} edge from parent node[above] {$\tfrac12$}}
   child {node {$T$} edge from parent node[below] {$\dots$}}
   edge from parent node[below] {$\dots$}};
\end{tikzpicture}
\end{center}
\begin{enumerate}
\item Compléter l'arbre en justifiant chaque calcul.
\item Calculer $P(E\cap R)$ et $P(\overline E\cap S)$.
\item Calculer $P(T)$.
\end{enumerate}
}{%
\begin{enumerate}
\item Les branches issues d'un même nœud ont pour somme $1$ : $P(\overline E)=1-\dfrac49=\dfrac59$ ; $P_E(S)=1-\dfrac14-\dfrac13=\dfrac5{12}$ ; $P_{\overline E}(T)=1-\dfrac3{10}-\dfrac12=\dfrac15$.
\item $P(E\cap R)=\dfrac49\times\dfrac14=\dfrac19$ et $P(\overline E\cap S)=\dfrac59\times\dfrac12=\dfrac5{18}$.
\item $E$ et $\overline E$ forment une partition, donc
$P(T)=P(E\cap T)+P(\overline E\cap T)=\dfrac49\times\dfrac13+\dfrac59\times\dfrac15=\dfrac4{27}+\dfrac19=\dfrac4{27}+\dfrac3{27}=\dfrac7{27}$.
\end{enumerate}
}{5}

\mxexdef{41}{Traduire un événement, produit sur un chemin}{%
Dans une salle de sport, $\pc{65}$ des inscrits ont un abonnement annuel, les autres un abonnement mensuel. Chaque inscrit choisit une activité principale. Parmi les abonnés annuels, $\pc{50}$ choisissent la musculation, $\pc{30}$ les cours collectifs et les autres la natation. Parmi les abonnés mensuels, $\pc{40}$ choisissent la musculation, $\pc{45}$ les cours collectifs et les autres la natation.

On choisit un inscrit au hasard. On note $A$ : « L'inscrit a un abonnement annuel », $M$ : « Son activité principale est la musculation », $C$ : « Son activité principale est un cours collectif » et $N$ : « Son activité principale est la natation ».
\begin{enumerate}
\item Construire un arbre pondéré décrivant la situation.
\item Décrire par une phrase l'événement $\overline A\cap C$, puis calculer sa probabilité.
\item Calculer la probabilité que l'inscrit choisi ait la natation pour activité principale.
\item L'inscrit choisi pratique principalement la natation. Calculer la probabilité qu'il ait un abonnement annuel. Arrondir à $10^{-3}$.
\end{enumerate}
}{%
\begin{enumerate}
\item Premier niveau : $P(A)=\num{0,65}$ et $P(\overline A)=\num{0,35}$. Sous $A$ : $P_A(M)=\num{0,50}$, $P_A(C)=\num{0,30}$, $P_A(N)=1-\num{0,50}-\num{0,30}=\num{0,20}$. Sous $\overline A$ : $P_{\overline A}(M)=\num{0,40}$, $P_{\overline A}(C)=\num{0,45}$, $P_{\overline A}(N)=1-\num{0,40}-\num{0,45}=\num{0,15}$.
\item $\overline A\cap C$ : « L'inscrit a un abonnement mensuel et son activité principale est un cours collectif ». $P(\overline A\cap C)=\num{0,35}\times\num{0,45}=\num{0,1575}$.
\item Par les probabilités totales ($A$, $\overline A$ forment une partition) : $P(N)=\num{0,65}\times\num{0,20}+\num{0,35}\times\num{0,15}=\num{0,13}+\num{0,0525}=\num{0,1825}$.
\item $P_N(A)=\dfrac{P(A\cap N)}{P(N)}=\dfrac{\num{0,13}}{\num{0,1825}}\approx\num{0,712}$.
\end{enumerate}
}{6}

\mxexdef{42}{Utiliser et tester l'indépendance}{%
Dans un centre d'appels, on choisit un appel au hasard. On note $A$ l'événement « Le client attend plus de $5$ minutes avant d'être pris en charge » et $B$ l'événement « La communication est coupée à cause d'un incident réseau ». On sait que $P(A)=\num{0,12}$ et $P(B)=\num{0,05}$. On suppose dans les questions 1 à 3 que $A$ et $B$ sont indépendants.
\begin{enumerate}
\item Calculer la probabilité que le client attende plus de $5$ minutes et que la communication soit coupée.
\item Calculer la probabilité qu'aucun de ces deux incidents ne se produise.
\item En déduire la probabilité qu'au moins un des deux incidents se produise.
\item Un audit montre en réalité que, parmi les appels avec une attente de plus de $5$ minutes, $\pc{9}$ sont coupés. L'hypothèse d'indépendance est-elle confirmée ? Justifier.
\end{enumerate}
}{%
\begin{enumerate}
\item Par indépendance, $P(A\cap B)=P(A)\times P(B)=\num{0,12}\times\num{0,05}=\num{0,006}$.
\item Si $A$ et $B$ sont indépendants, $\overline A$ et $\overline B$ le sont aussi : $P(\overline A\cap\overline B)=\num{0,88}\times\num{0,95}=\num{0,836}$.
\item L'événement « au moins un incident » est le contraire de « aucun incident » : sa probabilité est $1-\num{0,836}=\num{0,164}$.
\item L'audit donne $P_A(B)=\num{0,09}$, alors que $P(B)=\num{0,05}$. Comme $P_A(B)\neq P(B)$, $A$ et $B$ ne sont pas indépendants : une longue attente s'accompagne plus souvent d'une coupure.
\end{enumerate}
}{5}

\mxexdef{43}{Du tableau d'effectifs aux arbres pondérés}{%
Un groupe emploie $\num{1600}$ salariés répartis sur trois sites : le siège, l'usine et l'entrepôt. On a recensé les salariés pratiquant le télétravail au moins un jour par semaine.
\[
\begin{array}{l|ccc}
&\text{Télétravail}&\text{Pas de télétravail}&\text{Total}\\\hline
\text{Siège}&300&\cdots&400\\
\text{Usine}&\cdots&\cdots&\cdots\\
\text{Entrepôt}&\cdots&340&400\\\hline
\text{Total}&420&\cdots&\num{1600}
\end{array}
\]
\begin{enumerate}
\item Compléter le tableau.
\item On choisit un salarié au hasard. On note $S_1$ : « Il travaille au siège », $S_2$ : « Il travaille à l'usine », $S_3$ : « Il travaille à l'entrepôt » et $T$ : « Il pratique le télétravail ». Construire l'arbre pondéré dont le premier niveau est $S_1$, $S_2$, $S_3$ et le second $T$, $\overline T$.
\item Construire l'arbre pondéré dont le premier niveau est $T$, $\overline T$ et le second $S_1$, $S_2$, $S_3$. Donner les probabilités sous forme de fractions irréductibles, puis arrondies à $10^{-3}$.
\end{enumerate}
}{%
\begin{enumerate}
\item Siège : $400-300=100$ sans télétravail. Entrepôt : $400-340=60$ en télétravail. Usine : $\num{1600}-400-400=800$ salariés, dont $420-300-60=60$ en télétravail et $800-60=740$ sans. Total sans télétravail : $100+740+340=\num{1180}$ ($=\num{1600}-420$).
\item $P(S_1)=\dfrac{400}{\num{1600}}=\num{0,25}$, $P(S_2)=\dfrac{800}{\num{1600}}=\num{0,5}$, $P(S_3)=\num{0,25}$.\\
$P_{S_1}(T)=\dfrac{300}{400}=\num{0,75}$ et $P_{S_1}(\overline T)=\num{0,25}$ ; $P_{S_2}(T)=\dfrac{60}{800}=\num{0,075}$ et $P_{S_2}(\overline T)=\num{0,925}$ ; $P_{S_3}(T)=\dfrac{60}{400}=\num{0,15}$ et $P_{S_3}(\overline T)=\num{0,85}$.
\item $P(T)=\dfrac{420}{\num{1600}}=\dfrac{21}{80}=\num{0,2625}$ et $P(\overline T)=\dfrac{59}{80}=\num{0,7375}$.\\
$P_T(S_1)=\dfrac{300}{420}=\dfrac57\approx\num{0,714}$, $P_T(S_2)=\dfrac{60}{420}=\dfrac17\approx\num{0,143}$, $P_T(S_3)=\dfrac17\approx\num{0,143}$.\\
$P_{\overline T}(S_1)=\dfrac{100}{\num{1180}}=\dfrac5{59}\approx\num{0,085}$, $P_{\overline T}(S_2)=\dfrac{740}{\num{1180}}=\dfrac{37}{59}\approx\num{0,627}$, $P_{\overline T}(S_3)=\dfrac{340}{\num{1180}}=\dfrac{17}{59}\approx\num{0,288}$.\\
Contrôle : les fractions issues de chaque nœud ont pour somme $1$.
\end{enumerate}
}{7}

\mxexdef{44}{Inverser un conditionnement}{%
Une enseigne de bricolage interroge ses clients. $\pc{45}$ d'entre eux ont téléchargé l'application de fidélité de l'enseigne ; parmi ceux-ci, $\pc{60}$ ont effectué au moins un achat en promotion au cours du dernier mois. Au total, $\pc{36}$ des clients interrogés ont effectué au moins un achat en promotion au cours du dernier mois.

On choisit au hasard un client interrogé. On note $A$ : « Le client a téléchargé l'application » et $R$ : « Le client a effectué au moins un achat en promotion le mois dernier ».
\begin{enumerate}
\item Donner $P(A)$, $P(R)$ et $P_A(R)$.
\item Calculer $P(A\cap R)$, puis $P_R(A)$. Interpréter ce dernier résultat.
\item Calculer la probabilité qu'un client n'ayant pas téléchargé l'application ait effectué un achat en promotion. Arrondir à $10^{-3}$ et comparer avec $P_A(R)$.
\end{enumerate}
}{%
\begin{enumerate}
\item $P(A)=\num{0,45}$, $P(R)=\num{0,36}$ et $P_A(R)=\num{0,60}$.
\item $P(A\cap R)=P(A)\times P_A(R)=\num{0,45}\times\num{0,60}=\num{0,27}$. Puis $P_R(A)=\dfrac{P(A\cap R)}{P(R)}=\dfrac{\num{0,27}}{\num{0,36}}=\num{0,75}$ : les trois quarts des clients ayant acheté en promotion ont téléchargé l'application.
\item $P(\overline A\cap R)=P(R)-P(A\cap R)=\num{0,36}-\num{0,27}=\num{0,09}$ et $P(\overline A)=\num{0,55}$, donc $P_{\overline A}(R)=\dfrac{\num{0,09}}{\num{0,55}}\approx\num{0,164}$. C'est nettement moins que $P_A(R)=\num{0,60}$ : les clients ayant l'application achètent bien plus souvent en promotion.
\end{enumerate}
}{4}

\mxexdef{45}{Expérience en deux étapes, résultats fractionnaires}{%
Dans cet exercice, les résultats sont donnés sous forme de fractions irréductibles. Lors d'une animation commerciale, un client tire au hasard un ticket parmi dix tickets numérotés de $1$ à $10$.
\begin{itemize}
\item Si le numéro est $1$ ou $2$, il gratte une carte prise au hasard dans le coffret X, qui contient $3$ cartes gagnantes sur $8$.
\item Si le numéro est un multiple de $3$, il gratte une carte du coffret Y, qui contient $2$ cartes gagnantes sur $9$.
\item Sinon, il gratte une carte du coffret Z, qui contient $1$ carte gagnante sur $6$.
\end{itemize}
On note $X$, $Y$, $Z$ les événements « La carte provient du coffret X » (respectivement Y, Z) et $G$ l'événement « La carte est gagnante ».
\begin{enumerate}
\item Justifier que $P(X)=\dfrac15$, $P(Y)=\dfrac3{10}$ et $P(Z)=\dfrac12$, puis représenter la situation par un arbre pondéré.
\item Démontrer que la probabilité que la carte soit gagnante est égale à $\dfrac9{40}$.
\item Le client a gagné. Quelle est la probabilité que sa carte provienne du coffret X ?
\end{enumerate}
}{%
\begin{enumerate}
\item Les dix numéros sont équiprobables. $X$ correspond à $\ens{1;2}$ : $P(X)=\dfrac2{10}=\dfrac15$. $Y$ correspond à $\ens{3;6;9}$ : $P(Y)=\dfrac3{10}$. $Z$ correspond aux cinq autres numéros $\ens{4;5;7;8;10}$ : $P(Z)=\dfrac5{10}=\dfrac12$.\\
Arbre : sous $X$, $P_X(G)=\dfrac38$ et $P_X(\overline G)=\dfrac58$ ; sous $Y$, $P_Y(G)=\dfrac29$ et $P_Y(\overline G)=\dfrac79$ ; sous $Z$, $P_Z(G)=\dfrac16$ et $P_Z(\overline G)=\dfrac56$.
\item $X$, $Y$, $Z$ forment une partition :
$P(G)=\dfrac15\times\dfrac38+\dfrac3{10}\times\dfrac29+\dfrac12\times\dfrac16=\dfrac3{40}+\dfrac1{15}+\dfrac1{12}=\dfrac{9}{120}+\dfrac{8}{120}+\dfrac{10}{120}=\dfrac{27}{120}=\dfrac9{40}$.
\item $P_G(X)=\dfrac{P(X\cap G)}{P(G)}=\dfrac{3/40}{9/40}=\dfrac13$.
\end{enumerate}
}{6}

\mxexdef{46}{Bilan : chemin, contraire et répétition indépendante}{%
Un incubateur d'entreprises sélectionne les projets qu'il accompagne en trois étapes. $\pc{30}$ des dossiers déposés passent la présélection. Les projets présélectionnés sont présentés devant un jury, qui en retient $\pc{50}$. Enfin, un comité d'investisseurs finance $\pc{20}$ des projets retenus par le jury. Un projet éliminé à une étape ne passe pas les étapes suivantes.

On choisit un dossier au hasard. On note $P$ : « Le dossier passe la présélection », $J$ : « Le projet est retenu par le jury » et $F$ : « Le projet est financé ».
\begin{enumerate}
\item Construire un arbre pondéré décrivant la situation. Démontrer que la probabilité qu'un projet ne soit pas financé est égale à $\num{0,97}$.
\item Huit porteurs de projet déposent un dossier ; on admet que les dossiers sont étudiés indépendamment les uns des autres. Calculer la probabilité qu'au moins un des huit projets soit financé. Arrondir à $10^{-3}$.
\item Combien de dossiers l'incubateur doit-il recevoir, au minimum, pour que la probabilité de financer au moins un projet dépasse $\num{0,99}$ ?
\end{enumerate}
}{%
\begin{enumerate}
\item Arbre : $P$ ($\num{0,30}$) ou $\overline P$ ($\num{0,70}$, fin) ; après $P$ : $J$ ($\num{0,50}$) ou $\overline J$ ($\num{0,50}$, fin) ; après $J$ : $F$ ($\num{0,20}$) ou $\overline F$ ($\num{0,80}$). Un projet n'est financé que s'il franchit les trois étapes, donc, en multipliant le long de ce chemin :
$P(F)=\num{0,30}\times\num{0,50}\times\num{0,20}=\num{0,03}$ et $P(\overline F)=1-\num{0,03}=\num{0,97}$.
\item Par indépendance, la probabilité qu'aucun des huit projets ne soit financé est $\num{0,97}^8$. La probabilité cherchée est $1-\num{0,97}^8\approx\num{0,216}$.
\item On cherche le plus petit entier $n$ tel que $1-\num{0,97}^n>\num{0,99}$, c'est-à-dire $\num{0,97}^n<\num{0,01}$. La fonction $\log$ étant croissante, cela équivaut à $n\log(\num{0,97})<\log(\num{0,01})$, soit $n>\dfrac{\log(\num{0,01})}{\log(\num{0,97})}\approx\num{151,2}$ (on divise par $\log(\num{0,97})<0$, ce qui change le sens de l'inégalité). Il faut au minimum $152$ dossiers (contrôle : $1-\num{0,97}^{151}\approx\num{0,98994}$ et $1-\num{0,97}^{152}\approx\num{0,99024}$).
\end{enumerate}
}{6}
